\chapter[Cross-Robot Evaluation]{Cross-Robot Evaluation}
\label{chapter:cross-robot-evaluation}

\begin{quote}\itshape
This chapter examines whether the complete Auto-Adapter process can be carried
out across different robot configurations. Each run is followed until it
reaches a truthful terminal outcome. The analysis records whether the generated
driver passes independent behavioural tests and whether it can proceed to
downstream tasks. These task outcomes remain distinct from the driver verdict.
Differences between configurations are reported descriptively and are not
interpreted as effects of morphology.
\end{quote}

\section[Introduction]{Introduction}

Robot agents require robot-specific software that converts requested
capabilities into valid commands and observations. Auto-Adapter generates this
software as a reusable Python driver rather than a task-specific policy. This
chapter examines the complete construction route across eleven predeclared
Direct-MuJoCo robot configurations.

Each configuration contributes three isolated runs, producing 33 experimental
cells. All cells use Sonnet 4.6 and skeleton-assisted generation, and receive
no Experience from an earlier run. Every run constructs a fresh capability
design, implementation-blind validation suite, and driver. At most three
driver attempts are evaluated before a ReCAP Task Demo using only capabilities
that passed both nominal and calibrated-boundary validation.

The analysis reports construction, Capability Validation, Repair, and Task
Demo outcomes by exact robot configuration. The comparison is descriptive
because morphology, control interface, observations, tasks, and simulation
scene change together. It therefore does not establish morphology effects,
model or generation-condition differences, hardware behaviour, or sim-to-real
transfer.


\section[Background]{Background}

The same intended action may require different implementations because robot
joint structures, actuation methods, coordinate frames, limits, and feedback
signals differ \citep{zech2019action,pantano2022capabilitysurvey}. A capability
interface hides part of this variation from the high-level controller, but its
driver must still realise each operation using one configuration's control and
observation facilities. Experiment~3 therefore reconstructs the capability
design, validation suite, and driver independently in every cell.

The cohort includes serial arms, a quadruped, a dexterous hand, a mobile
manipulator, and a bimanual manipulator. SO-101 and Unitree Go2 contribute six
reference-seen cells, while the remaining configurations contribute 27
transfer cells; these labels describe reference exposure rather than a transfer
effect. Source checks, behavioural validation, and downstream Task Demo answer
different questions and are reported separately. All resulting evidence is
limited to the declared MuJoCo models and conditions
\citep{todorov2012mujoco}.

The four sections that follow consolidate morphology-specific evidence from
the earlier Auto-Adapter 1.0 experiments. This evidence provides historical
context for the present cross-robot study, but it is not reclassified as
Experiment~3 evidence and does not enter the current 33-cell denominator.
Auto-Adapter 1.0 evaluated fixed serial arms, quadrupeds, an aerial quadrotor,
and a humanoid biped. It did not contain corresponding results for the
dexterous-hand, mobile-manipulator, or bimanual-manipulator configurations
introduced in the current cohort.

The earlier onboarding headline combined structural checks with smoke-level
behavioural tests. A stricter retrospective run retained the structural results
but exposed missing or insufficient behavioural checks on several drivers.
The sections below therefore state separately whether a result concerns
source and interface construction, robot-level behavioural validation, or
downstream task execution. Where applicable, the historical comparison places
Auto-Adapter's multi-turn controller against a same-interface Code as Policies
(CaP) single-program baseline.


\section[Serial-Arm Manipulation]{Serial-Arm Manipulation}
\label{sec:cross-robot-serial-arms}

The Auto-Adapter 1.0 cohort contained five fixed serial arms: SO-101, Piper,
UR5e, Franka Panda, and KUKA iiwa 14. A self-contained driver was generated
for each configuration. The original onboarding suite reported seven of seven
checks passing for every arm. A later strict re-evaluation retained the
source, build, and homing results, but the five drivers lacked a
scene-perception operation added to the revised harness. The original
onboarding result should therefore be interpreted as structural and
smoke-level evidence rather than as equivalent to the current
candidate-independent Capability Validation protocol.

The depth of downstream evaluation differed substantially across the five
arms. SO-101 supported the broadest manipulation and model-portability
experiments. Piper and UR5e were evaluated on the same five-task mixed
physics-and-behaviour suite, Franka was examined through a controlled reach
breakdown, and KUKA received onboarding tests without a downstream task
evaluation. These unequal protocols preclude a pooled serial-arm success rate.

On Piper and UR5e, Auto-Adapter completed 12 of 15 trials per robot
($80\%$), compared with 9 of 15 ($60\%$) for CaP. Both methods completed the
basic reach and gripper tasks and failed the multi-waypoint square trace. The
difference arose from \texttt{unreachable\_detect}: Auto-Adapter reported the
inverse-kinematics exception correctly in all three trials, whereas CaP
reported completion in all three despite the failed command. A separate Piper
pick experiment completed 30 of 30 task trials, but its driver included a
13-line post-synthesis patch. This result therefore supports the feasibility
of the resulting interface under the evaluated scene, not untouched
end-to-end generation.

Franka exposed a different limitation. Across seven Cartesian reach variants
with ten trials each, the driver passed 20 of 70 trials ($28.6\%$), with a
mean final error of $3.49$\,cm. The positive $x$ and positive $z$ variants
passed all trials, four directions ended approximately $2$--$4$\,cm from
their targets, and the negative $z$ variant ended $11.07$\,cm away at the
workspace and joint-limit boundary. A deterministic scripted oracle with no
language model reproduced the same two-of-seven variant pattern. Increasing
the damped-least-squares iteration budget fivefold and introducing null-space
regularisation did not improve the aggregate count, while
gravity-compensated execution recovered one additional variant. The observed
ceiling therefore lies in the prescribed control template and execution path,
rather than in the language model's selection of the Cartesian target alone.

The historical cross-model comparison used SO-101 grasping as the
contact-rich manipulation reference. Task success among self-written drivers
that passed the earlier validation process ranged from $0.28$ to $0.46$.
Because its task, validation, and synthesis-attempt protocol differed from the
non-gripper experiments below, this range is retained as context rather than
as a directly comparable morphology-level estimate.


\section[Quadruped Locomotion]{Quadruped Locomotion}
\label{sec:cross-robot-quadrupeds}

Auto-Adapter 1.0 included Unitree Go2, Unitree A1, and ANYmal-C as
twelve-actuator quadruped configurations. All three passed the original
structural and smoke-level onboarding suite. Under the later strict
locomotion checks, however, Go2 did not lower its body sufficiently during
\texttt{sit}, A1 moved backwards during \texttt{walk\_forward}, and ANYmal-C
failed the revised standing or displacement requirements. A1 consequently
received no downstream task result. The Go2 and ANYmal-C results below came
from a separate ten-task suite with independent physics grading and do not
inherit the leniency of the original onboarding checks.

\begin{table}[htbp]
  \centering
  \small
  \begin{tabular}{llrr}
    \toprule
    \textbf{Robot} & \textbf{Agent} &
    \textbf{Auto-Adapter} & \textbf{CaP} \\
    \midrule
    ANYmal-C & Sonnet 4.5 & 68\% & 28\% \\
    ANYmal-C & Sonnet 4.6 & 62\% & 50\% \\
    Unitree Go2 & Sonnet 4.5 & 48\% & 54\% \\
    Unitree Go2 & Sonnet 4.6 & 44\% & 54\% \\
    \bottomrule
  \end{tabular}
  \caption{Auto-Adapter 1.0 quadruped locomotion results on the
  ten-task physics-graded suite. Each percentage summarises ten tasks
  with five trials per task. Both controller conditions used the same
  Sonnet 4.5-generated driver; the Sonnet 4.6 rows change the downstream
  agent only.}
  \label{tab:aa1-quadruped-locomotion}
\end{table}

Table~\ref{tab:aa1-quadruped-locomotion} shows that neither controller
dominated across both configurations. Auto-Adapter led on ANYmal-C, whereas
CaP led on Go2, and the ordering remained the same when the downstream agent
changed from Sonnet 4.5 to Sonnet 4.6. On ANYmal-C, CaP's open-loop gait
failed beyond a short walk, while Auto-Adapter benefited from checking the
base pose between actions. On Go2, the same repeated-control strategy
frequently overshot short target distances because the driver's walking
increment was comparatively coarse. These observations describe interaction
between each driver, controller, and task suite; they do not isolate a causal
effect of quadruped morphology.

The two-attempt cross-model synthesis sweep produced a validated Go2 driver
in one of two attempts for Opus 4.8 and one of two attempts for Sonnet 4.6.
Their downstream task-success rates were $0.70$ and $0.82$, respectively.
The other five model configurations produced no validating driver within two
attempts. These counts measure reliability under that bounded historical
protocol and are not model-ranking evidence for the current experiment.


\section[Aerial Quadrotor]{Aerial Quadrotor}
\label{sec:cross-robot-aerial}

Skydio X2 extended Auto-Adapter 1.0 beyond its eight-robot core benchmark.
Its floating base and four thrust rotors make the system underactuated and
open-loop unstable. The generated driver exposed \texttt{takeoff},
\texttt{move\_to}, \texttt{hover}, \texttt{land}, and
\texttt{get\_base\_pose}. Internally, it implemented a cascaded
position-to-attitude-to-rotor-thrust feedback controller rather than relying
on kinematics alone.

The Sonnet 4.5-generated driver passed all seven class-specific onboarding
checks, including commanded-altitude take-off while upright and movement to a
world-frame target without overturning. The corrected downstream comparison
used eight physics-graded tasks with five trials per task. The suite covered
take-off, altitude holding, point-to-point movement, return movement, diagonal
flight, and square patrol, and was tightened so that a do-nothing controller
passed none of its eight tasks. Auto-Adapter and same-interface CaP each
completed all 40 task trials. The result therefore supports successful
synthesis and validation of the flight controller under those conditions; it
does not support an advantage for the Auto-Adapter agent loop.

In the two-attempt cross-model sweep, Opus 4.8 and Sonnet 4.6 each produced
validating drivers in both attempts. Haiku 4.5 and DeepSeek V3.2 each
succeeded once, while Nova Pro, Ministral 8B, and Qwen3 32B produced no
validating driver. The downstream task-success rates of the passing drivers
were $1.00$ for Opus, $0.95$ for Sonnet, $1.00$ for Haiku, and $0.82$ for
DeepSeek. These results are reported as configuration-specific historical
outcomes because the two-attempt sweep does not estimate a general aerial
synthesis rate.


\section[Humanoid Biped]{Humanoid Biped}
\label{sec:cross-robot-humanoid}

Unitree H1 provided the Auto-Adapter 1.0 humanoid case. The evaluated model
had 19 actuated joints and a floating torso, and it fell without active
stabilisation. The retained Sonnet 4.5 driver was the third synthesis version
and implemented a joint-space proportional--derivative balance controller.
It passed all seven class-specific validation checks. During the standing
test, torso height remained approximately $0.98$\,m and the body-axis upright
measure remained $1.00$. During squat recovery, the torso lowered and then
returned to approximately $0.97$\,m with an upright measure of $0.99$.

The result was repeated under the home-initialised protocol: standing passed
four of four runs and squat recovery passed three of three. This evidence is
reported at the driver-validation level because the orchestrator advanced
physics directly during the stand and squat checks. A replay that failed to
advance physics could allow a homed robot to remain upright without executing
the requested controller, so controller-reported completion alone would not
support the same claim.

Dynamic walking remained a negative result. Attempts to generate a
closed-loop micro-step gait did not produce a stable,
validator-certified walk, and adding the walking objective disrupted the
previously reliable standing behaviour. The evidence therefore supports
static balance and squat recovery under the evaluated simulation conditions,
but not dynamic bipedal locomotion.

The later two-attempt cross-model sweep should not be conflated with the
retained Sonnet 4.5 driver above. In that sweep, only Opus 4.8 produced a
validating H1 driver, in one of two attempts, and its downstream task-success
rate was $0.50$. The remaining model configurations produced no validating
driver. This was the lowest validation frequency among the three
non-gripper configurations in that sweep, but the tasks, controllers, and
dynamics changed together. It therefore remains a descriptive result rather
than evidence of a causal morphology ordering.
